\section{Related Work}\label{sec:related}

\subsection{Verification Problems and Their Significance}

This paper focuses on verification problems for neural networks and, in particular, the range analysis problem.  Neural networks are now used in safety-critical applications such as self-driving cars \cite{bojarski_testa_2016}, robotics \cite{gu2017deep}, and smart prosthetics \cite{clark_campbell_amor}. In this context, they suffer from some unique challenges beyond incorrect decisions arising from lack of training data and over-fitting. Adversarial attacks were first reported by Goodfellow et al.~\cite{goodfellow_shlens_szegedy_2015}, where small input perturbations on datasets like MNIST-10 (handwritten digits) led to wrongful classifications with high probability. Since then, several techniques have been developed to train models to prevent such attacks. A detailed survey can be found here \cite{bai2021recent,silva2020opportunities}.

Verification tools have  addressed these issues by proving that a given neural network is robust to certain types of perturbations or synthesizing a counterexample that shows a vulnerability in the network. Existing tools have been focused on feedforward neural networks. They  include Reluplex~\cite{katz_barrett_dill_julian_kochenderfer_2017}, \textsc{Marabou} (SMT-based) \cite{katz_huang_2019}, \textsc{CROWN} (uses convex relaxations) \cite{zhang_weng_chen_hsieh_daniel_2018},  \textsc{Sherlock} (MILP-based) \cite{dutta_chen_jha_sriram_sankaranarayanan_tiwari_2019} and others~\cite{lomuscio_maganti_2017} (see~\cite{Huang2020Survey,Liu+Others/2021/Algorithms,albarghouthi-book} for a detailed survey). These tools reduce the complex verification question to  range verification, the precise problem being tackled in this paper. 
They have been used to certify safety of an aircraft collision avoidance system~\cite{Irfan+Others/2020/Towards,katz_barrett_dill_julian_kochenderfer_2017}, prove bounds on inputs to control systems~\cite{Majid+Others/2022/Safe}, study the sensitivity of cyber-physical models to perturbations~\cite{Kushner+Sankaranarayanan+Breton/2020/Conformance} and certify robustness to certain types of adversarial attacks~\cite{tjeng_xiao_tedrake_2017,Singh+Others/2019/Beyond}. 
%As we explore the applications to KANs in these domains, the same verification problems need to be solved for KANs, as well. 

\subsection{NNs with Learnable Activation Functions}

Neural networks with  learnable ``free form'' activation functions such as deep-spline networks have been shown to be advantageous over those with fixed activation functions \cite{Aziznejad+Others/2020/Deep,Pakshal+Others/2020/Learning} in terms of being able to learn more succinct and accurate function representations but at the cost of increased training time. More recently,  Kolmogorov Arnold Networks (KANs) also replace fixed nonlinear activation functions with general univariate functions that can change during the training process. This leads to networks that are more parameter-efficient than MLPs, while still maintaining competitive or superior performance~\cite{Liu+Others/2025/KAN}.  KANs have been shown to achieve better accuracies when compared with larger MLPs for the same learning task~\cite{howard_jacob_murphy_heinlein_stinis_2024,toscano_wang_karniadakis_2024}. Applications  include time-series predictions~\cite{genet2024tkan}, vision \cite{mahara2024dawn} and physics-based learning \cite{abueidda2025deepokan}. Due to the nonlinear and free-form nature of these activation functions, these networks are not amenable to existing range verification approaches beyond those based on simple interval bound propagation~\cite{moore2009introduction}, which suffers from the well-known wrapping effect and consequently tends to yield very coarse bounds on networks of this kind.

\subsection{Neural Network Abstraction}

 Abstracting a neural network with activation functions such as tanh, sigmoid and gelu by means of a PWA function has been investigated as part of most range verification approaches~\cite{albarghouthi-book}. 
 Pulina and Tacchella  were one of the first to propose a form of piecewise affine abstractions by replacing each nonlinear activation function by a series of intervals~\cite{Pulina+Tacchella/2012/Challenging}. Their abstraction of nonlinear activation functions subdivides the domain uniformly using a single width parameter $p$ and uses the fact that commonly used activation functions (sigmoid and tanh) are monotonically increasing functions to bound the value of the function in each subinterval. If a given large value of the width parameter is unable to prove the property, the refinement process lowers the width produce a finer grained abstraction.  In contrast, our work here uses a PWA approximation that is optimized in two senses: we use Bellman's algorithm at each node to find the best way to grid the domain into $k$ pieces that minimizes the approximation error~\cite{Bellman/1961/Approximation,Bellman+Roth/1969/Curve}, and we use different number of pieces at different nodes. While refinement is not part of the approach in this paper, we can implement a refinement loop that is similar in spirit to the work of Pulina and Tacchella by incrementally increasing the total number of pieces used in our abstraction. 
 
 Following Pulina and Tacchella, most subsequent approaches either focus exclusively on networks with piecewise affine activation functions such as ReLU~\cite{lomuscio_maganti_2017,katz_barrett_dill_julian_kochenderfer_2017,Ehlers/2017/Formal,Tran+Othrs/2019/Star}, or employ an \emph{a priori} fixed abstraction that is performed once for each activation function to achieve a desired local error bound (see for example~\cite{dutta_chen_jha_sriram_sankaranarayanan_tiwari_2019}). Accounting for this error yields a sound output range that can be used to establish properties. While such approaches can work for activation functions such as sigmoid and tanh that are well approximated by PWA functions with up to $5$ pieces, they fail for networks that employ B-splines or radial basis functions as their activation functions. 

 Bound propagation approaches such as CROWN~\cite{zhang_weng_chen_hsieh_daniel_2018} or DeepPoly~\cite{Singh+Others/2019/Beyond} can eliminate the need for an upfront abstraction by generating upper and lower bounds for nonlinear activation functions on-the-fly as the analysis is performed. At the same time, such approaches can yield coarse bounds if the input bounds are large and the activation function is not monotone. The recent work of  Shi et al.~\cite{Shi+Others/2025/Neural} improves the handling nonlinear activation functions such as tanh, sinh and variants of ReLU using a generalized branch-and-bound approach that iteratively branches nonlinear activation functions at carefully chosen points, and uses the bound propagation approach of alpha-beta-CROWN~\cite{Wang+Others/2021/BetaCROWN}. Shi et al.\ assumes a pre-optimized set of branching points for their approach for a fixed library of commonly occurring nonlinear functions and chooses a single branching point on the fly, whereas our approach computes PWA approximation upfront for a given error tolerance $\delta$ through dynamic programming. In a sense, the contributions of our work can likely improve the approach of Shi et al.\ by providing a principled approach for choosing branching points.

 The work of Ivanov et al.\ uses a creative approach that encodes neural networks as nonlinear differential equations and converts neural network verification problems to hybrid systems verification problems~\cite{IvanovEtAl2019Verisig}. While activation functions such as tanh and sigmoid can indeed be seen as satisfying a differential equation, it is not clear how arbitrary activation functions can be encoded in this framework. 

\subsection{Verification for KANs and B-Spline Networks}
To the best of our knowledge, few publications have explored the verification of KANs, and the field remains relatively unexplored. The related problem of robustness of  KANs to adversarial perturbation has been investigated: Heiderich et al.~\cite{heiderich2025training} present a comparative analysis of KANs and MLPs in the context of adversarial robustness, concluding that standard KAN architectures are similarly vulnerable and introducing a new training method.
\cite{polomolina2024monokancertifiedmonotonickolmogorovarnold} modify the KAN architecture by using Cubic Hermite splines in place of B-Splines to guarantee monotonicity: the range analysis problem for such networks is solved simply by evaluating the network at two corner-points of the input hyper-rectangle.  Additional efforts include the works of \cite{shen2025reduced} and \cite{zeng2024kan}, where they similarly highlight the importance of adversarial robustness and study techniques to improve the effectiveness of KANs against added noise. Our approach is the first to consider the problem of systematically computing PWA abstractions of KANs using a reduction to a variant of the knapsack problem. 

% The need to develop interpretable, expressive yet efficient neural network architectures have led to the development of Kolmogorov-Arnold Networks (KANs) \cite{Liu+Others/2025/KAN}. Differing from traditional Multi-Layer Perceptron (MLP) architecture, where networks primarily learn a linear transformation interspersed with fixed non-linear activations functions, KANs draw inspiration from Kolmogorov-Arnold representation theorem \cite{Kolmogorov/1957/Representation}, \cite{Arnold/1958/Representation}. The theorem posits that any continuous multivariate function can be expressed as a sum of univariate functions. Thereby, KANs feature learnable activation functions, parameterized as splines, while the nodes perform summation over the functions. This shift in the universal approximator shows us several benefits, including the ability to achieve comparable (better in some cases) accuracies to larger MLPs \cite{howard_jacob_murphy_heinlein_stinis_2024}, \cite{toscano_wang_karniadakis_2024}.
% \ck{
% The need to develop interpretable, expressive yet efficient neural network architectures have recently led to the development of Kolmogorov-Arnold Networks (KANs) \cite{Liu+Others/2025/KAN}. Differing from traditional Multi-Layer Perceptron (MLP) architecture, where networks primarily learn a linear transformation interspersed with fixed non-linear activations functions, KANs draw inspiration from Kolmogorov-Arnold representation theorem \cite{Kolmogorov/1957/Representation}, \cite{Arnold/1958/Representation}. The theorem posits that any continuous multivariate function can be expressed as a sum of univariate functions. Thereby, KANs feature learnable activation functions, parameterized as splines, while the nodes perform summation over the functions. This shift in the universal approximator shows us several benefits, including the ability to achieve comparable (better in some cases) accuracies to larger MLPs \cite{howard_jacob_murphy_heinlein_stinis_2024}, \cite{toscano_wang_karniadakis_2024}. However, the increased representational power and the dynamic nature of these spline-based activations also introduce new complexities, notably in training efficiency and, critically for this work, in formal verification. \\
% Formal verification of neural networks has become a increasingly vital area of research, especially since the deployment of such systems are now common in safety-critical settings such as autonomous vehicles and medical diagnosis \cite{katz_barrett_dill_julian_kochenderfer_2017}, \cite{matthias_könig_bosman_hoos_jan_2024}. This method aims to provide mathematical guarantees that a neural network satisfies key properties such as robustness to adversarial inputs, output range bounds, and fairness over specified input domains. Despite significant progress, verifying neural networks remains a challenging task due to their high dimensionality and non-linear operations \cite{boudardara_boussif_meyer_ghazel_2024}. Existing verification techniques, such as those based on Satisfiability Modulo Theories (SMT) solvers, abstract interpretation, or reachability analysis, often struggle with the scalability and precision required for deep networks with highly non-linear components \cite{singh_gehr_püschel_vechev_2019}. Mixed-Integer Linear Programming (MILP) has emerged as a powerful verification method, particularly for networks employing piecewise affine (PWA) activation functions like ReLU \cite{tjeng_xiao_tedrake_2017}. However, the direct application of MILP to KANs is not well defined due to the non-PWA nature of the spline activations. \\
% To bridge this gap, abstraction techniques are often employed. Abstraction involves creating a simpler, yet conservative, mathematical model of the network (or its components) that can be much easier to verify, while still providing guarantees about the original network. Piecewise Affine (PWA) functions are a common choice for such abstractions, as they can be effectively encoded as a MILP. \\
% The proposed method develops a multi-stage optimization process. Drawing inspiration from classic curve-fitting techniques \cite{Bellman/1961/Approximation}, \cite{Bellman+Roth/1969/Curve} dynamic programming algorithms are utilized to determine the optimal PWA abstraction for each individual univariate spline function in the KAN, considering various numbers of affine pieces and their associated approximation errors. This stage effectively generates a set of candidate abstractions for each KAN unit, each with a known error and complexity.
% }


% \noah{
%Since the original KAN work, several papers have discussed methods for robustness verification. \cite{polomolina2024monokancertifiedmonotonickolmogorovarnold} modified the KAN architecture by using cubic Hermite splines in place of B-Splines and positive weights along the edges of the network to achieve certified partial monotonicity. Importantly, the paper obtains guarantees entirely mathematically because of their network's unique structure. In contrast, we investigate swift MILP verification of the original KAN architecture. Similarly, \cite{heiderich2025training} also investigates augmenting the KAN architecture before verification. They document an efficient method for computing Lipschitz constants which modify the network's output to encourage robustness. Crucially, the authors then 

% \todo[inline]{
% 1. Talk about KANs and the significance of that line of research. 1 paragraph : 2-3 lines.

% 2. Talk about NN verification and some of the key contributions, tools. Mention the survey/books in the area, tools such as Marabou, CROWN, and mention the applications to robustness to perturbations. Also, pay special attention to MILP based approaches including Sherlock, the work of Lomuscio et al.\ and Tedrake et al.\ among others.

% 3. Talk about KAN verification approaches and in each instance try to differentiate how they compare with our work.

% }
% }